\documentclass[12pt]{article}
\usepackage{amssymb, amsmath, amsfonts, bm, graphicx, geometry}
\usepackage{tabularx, booktabs, float, enumitem, multicol, mathtools}
\usepackage{parskip} % Auto-manages paragraph spacing and removes indentations
\usepackage{siunitx}  % Eases typesetting of numbers and units

% US customary units
\DeclareSIUnit{\inch}{in}
\DeclareSIUnit{\foot}{ft}
\DeclareSIUnit{\yard}{yd}
\DeclareSIUnit{\mile}{mi}

\DeclareSIUnit{\slug}{slug}
\DeclareSIUnit{\lbm}{lbm}
\DeclareSIUnit{\lbf}{lbf}

\DeclareSIUnit{\psi}{psi}
\DeclareSIUnit{\psf}{psf}

\DeclareSIUnit{\gallon}{gal}
\DeclareSIUnit{\quart}{qt}
\DeclareSIUnit{\fahrenheit}{F}
\DeclareSIUnit{\rankine}{R}
\sisetup{
    per-mode = fraction
}
\usepackage{soul}
\usepackage{mathrsfs}

\geometry{margin=1in}

% Uppercase Volume with thicker, higher line
\newcommand{\vol}{\mathord{\text{\ooalign{\hidewidth $V$\hidewidth\cr\kern-0.5pt\raisebox{0.8ex}{\rule[0pt]{1em}{0.5pt}}}}}}

% Lowercase Volume with built-in shifts
\newcommand{\svol}{%
  \mathord{\text{\ooalign{%
    \hfil$v$\hfil\cr
    \hfil\rule[0.6ex]{0.8em}{0.5pt}\hfil
  }}}%
}

\begin{document}
\begin{center}
    {\LARGE\bfseries Chapter 4: Fluid Kinematics\par}
\end{center}
Velocity and acceleration fields. Material derivatives (derived from chain rule on spatially and time dependent fields). Reynolds transport theorem.

% ==========================================
% CHAPTER 4: FLUID KINEMATICS
% ==========================================

\section*{\fbox{4.2 The Acceleration Field}}
\textbf{Problem 4.8.27 (9th\#4.2.5: Acceleration Field)}\\
\textbf{Problem 4.8.34 (9th\#4.2.12: Acceleration Field)}\\
\textbf{Challenge Activity 4.2.2}\\

% ==========================================
% CHAPTER 5: FINITE CONTROL VOLUME ANALYSIS
% ==========================================
\begin{center}
    {\LARGE\bfseries Chapter 5: Finite Control Volume Analysis\par}
\end{center}
Application of the Reynolds transport theorem to conservation laws and laws of thermodynamics.
\section*{\fbox{5.1 Conservation of Mass}}
\textbf{Problem 5.9.1 (9th\#5.1.1: Cons of Mass)}\\
\textbf{Problem 5.9.3 (9th\#5.1.3: Fixed nondeforming)}\\

\textbf{Problem 5.9.38 (9th\#5.2: Linear Momentum)}\\
\textbf{Problem 5.9.45 (9th\#5.2.11: Linear Momentum)}\\
\textbf{Problem 5.9.83 (9th\#5.2.49: Moment of Mom.)}\\
\textbf{Problem 5.9.85 (9th\#5.2.51: Moment of Mom.)}\\
\textbf{Challenge Activity 5.2.1}\\
\textbf{Challenge Activity 5.2.2}\\

\textbf{Problem 5.9.99 (9th\#5.3.9: Energy)}\\
\textbf{Problem 5.9.107 (9th\#5.3.1: Energy)}\\

% ==========================================
% CHAPTER 6: DIFFERENTIAL ANALYSIS OF FLUID FLOW
% ==========================================
\begin{center}
    {\LARGE\bfseries Chapter 6: Differential Analysis of Fluid Flow\par}
\end{center}
Application of Navier stokes equations.
\textbf{Problem 6.14.12 (9th\#6.2.2: Cons. of Mass)}\\
\textbf{Problem 6.14.23 (9th\#6.3.1: Linear Momentum)}\\
\textbf{Problem 6.14.86 (9th\#6.9.16: Couette Flow)}\\
\textbf{Challenge Activity 6.9.1}\\




\end{document}
